\section{从数据口径到决策稳健性的跨问题检验}
\label{sec:cross-validation}

单问指标只说明局部模型是否可用，跨问题审计则检查上游口径、参数和不确定性传到下游后会不会改变经营结论。表\ref{tab:cross-checks}只保留影响任务链衔接的核验结果，预测误差、因素效应和层级数值的完整解释见对应问题章节。

\begin{longtable}{>{\centering\arraybackslash}p{0.13\textwidth}>{\centering\arraybackslash}p{0.23\textwidth}>{\centering\arraybackslash}p{0.29\textwidth}>{\centering\arraybackslash}p{0.23\textwidth}}
\caption{跨问题口径、参数与决策审计表}\label{tab:cross-checks}\\
\toprule
传递链路 & 审计设置 & 核验结果 & 下游约束 \\
\midrule
\endfirsthead
\multicolumn{4}{c}{续表~\thetable}\\
\toprule
传递链路 & 审计设置 & 核验结果 & 下游约束 \\
\midrule
\endhead
\bottomrule
\endfoot
治理到识别 & 疑似复制保留或删除；异常发货日期修正或置缺 & 总销售额变化 $0.0124\%$，履约效应量差 $0.000131$，显著性结论一致 & 销售结构沿用去重主口径，物流只作双口径关联解释 \\ % src:求解/问题1/结果/验证与灵敏度.json:朴素基线;求解/问题2/结果/物流双口径灵敏度.json:效应量绝对差,显著性结论一致
识别到策略 & 平滑、树复杂度与跨年判据构成 $81$ 组网格 & 三项产品因素均为 $81/81$ 入选，区域仅为 $54/81$ & 产品信号进入重点动作，区域只保留次级提示 \\ % src:求解/问题2/结果/关键参数灵敏度.json:组合总数,各因素入选次数
验证到应用 & 2025验证集复用；逐日在线特征；固定2024-12-31离线预测 & 2025年已用于时间外评估与灵敏度诊断，不再视为完全未触碰的外部测试 & 逐日识别可使用当日之前已实现记录；若在年末一次性预测下年全年，须将历史特征截止日固定为2024-12-31并重新验证 \\ % src:交接/结果解读_问题2.md:二.4防泄漏,五.2,五.4;求解/问题2/结果/时间外分组置换重要性.json:训练区间,测试区间
预测到部署 & 季节锚下限与复杂度惩罚形成 $9$ 组扰动 & 各设定均保持主指标 $3/3$、季度 $4/4$ 优于基线 & 采用训练期锁定组合，不按测试年重新挑权重 \\ % src:求解/问题3/结果/组合权重与灵敏度.json:关键参数灵敏度;交接/实验记录.json:[问题=问题3/权重约束灵敏度复核]
总量到层级 & 近期权重与收缩强度形成 $9$ 组扰动 & 份额WAPE为 $29.010\%$--$31.031\%$，低于固定份额的 $39.585\%$ & 收缩优势稳定，细层名次仍受总量情景约束 \\ % src:求解/问题4/结果/份额收缩回测与灵敏度.json:关键参数灵敏度_2025仅作检验不改锁定参数,2025五维平均对比
层级到预算 & 五维点预测与区间边界逐月回加 & 最大总量误差 $5.82\times10^{-11}$，产品父子误差 $1.46\times10^{-11}$ & 各业务预算共享同一总量，区间仅作可加总风险边界 \\ % src:求解/问题4/结果/2026总量与一致性.json:层级一致性检查,区间解释
\end{longtable}

我们最初以 $0.75$ 的跨年方向一致率门槛把区域列入核心因素，但把门槛提高到 $1.00$ 后，区域因实际一致率仅为 $0.75$ 而退出；完整网格又出现重要性方向翻转，因此将区域降为次级信号，避免不稳定地理标签直接放大为强投入建议。 % src:交接/实验记录.json:[问题=问题2/正确性复核的稳定阈值灵敏度],[问题=问题2复核/关键参数翻转与区域降级];求解/问题2/结果/跨年稳定性.json:结果[因素=区域]

模型的优点在于四问共享可追溯口径：订单簇控制明细相关性，时间外基线约束预测复杂度，总量锚与父子协调保证预算可加总；任何信号只有经口径、参数和时间检验后才提高策略强度。这种把预测精度与层级一致性同时纳入评价的做法，也符合时间序列部署和层级预测的基本要求\cite{hyndman2021forecasting,hyndman2011hierarchical}。

局限来自数据长度与字段边界：短月度序列难以充分校准远期区间，细层份额误差没有与总量误差建立联合分布；利润、成本、折扣、促销、销量和库存字段缺失，又使稳定关联不能转写为干预因果。模型适合滚动更新销售额与容量优先级，不适合承诺利润最优、价格弹性或安全库存件数。
